%% LaTeX2e class for student theses
%% sections/abstract_en.tex
%% 
%% Karlsruhe Institute of Technology
%% Institute of Information Security and Dependability
%% Software Design and Quality (SDQ)
%%
%% Dr.-Ing. Erik Burger
%% burger@kit.edu
%%
%% Version 1.6, 2024-06-07

\Abstract
The efficiency of energy storage via a metal carrier largely depends on the efficiency of both reducing iron oxide to iron and oxidizing iron for the release of energy.
For the reduction recent experiments have used a fluidized bed reactor. 
The fluid dynamics, particle motion and particle-particle interactions inside a fluidized bed are complex.
One approach to better understand the flow and behavior inside a reactor is to use numerical simulations.
To reduce computation time particle-particle interactions are modeled based on properties calculated on an Eulerian mesh instead of resolved, using the \MP method. 
%We discuss the fluidization regimes and that we want to stay above the minimum fluidization velocity and below the terminal velocity so gas can pass through the particles efficiently but particles do not leave the reactor.
We look in detail into the \MP as it was first presented along with submodels used, which includes the packing model that prevents over-packing by accelerating particles when particles are near close pack.
The discretization was based on a simple finite differencing scheme, which is applied to the governing equations of the fluid and the particle phase.
We reformulated the particle properties used to only depend on the values of the Eulerian mesh and ended up with four equations: the discretized continuity and momentum equations for both the fluid and the solid phase, which were solved using provisional particle positions.
This differs from the implementation in \OF.
The particles are tracked explicitly, therefore the particles are moved based on Newton's equations of motion.
Several submodels are included in \OF. Here we discuss the isotropy model and the explicit Harris and Crighton packing model as well as two drag models and the \lstinline|PIMPLE| algorithm which solves the needed properties for the submodels via a \lstinline|PIMPLE| algorithm.
We conduct a grid study and estimate an order of convergence.
For a validation of the solver we process data provided by the "Division of Combustion Technology".
We compare the resulting bed height and pressure drop to simulated results for the Harris and Crighton model combined with different drag models and an isotropy model in a parametric study.
We show that the \MP method for a good parameter choice can predict the bed height well but deviates from the experimentally predicted pressure drop.